\documentclass[conference]{IEEEtran}

\usepackage{graphicx}
\usepackage{amsmath,amssymb}
\usepackage{float}
\usepackage{booktabs}
\usepackage{multirow}

\begin{document}

\title{
    \textbf{Benchmark Report: 511.povray\_r} \\
    \text{Register Spill Analysis on RISC-V via gem5}
}

\author{
    \IEEEauthorblockN{Lütfullah Burak Kaya}
    \IEEEauthorblockA{
        Ozyegin University \\
        burak.kaya.50711@ozu.edu.tr
    }
    \and
    \IEEEauthorblockN{Ismail Akturk}
    \IEEEauthorblockA{
        Ozyegin University \\
        ismail.akturk@ozyegin.edu.tr
    }
}

\newcommand{\LongDate}{%
    \number\day\space
    \ifcase\month
    \or January\or February\or March\or April\or May\or June%
    \or July\or August\or September\or October\or November\or December%
    \fi
    \space \number\year
}

\twocolumn[
\begin{@twocolumnfalse}
\maketitle
\begin{center}{\small \LongDate}\end{center}
\vspace{1em}
\end{@twocolumnfalse}
]

% ================================================

\begin{abstract}
This report presents the register spill characterization of the
\texttt{511.povray\_r} benchmark from SPEC CPU2017, executed on a
RISC-V (RV64GC) target using the gem5 TimingSimpleCPU simulator.
A custom SpillDetector was integrated into gem5 to count store-then-load
patterns on the stack within a user-defined Region of Interest (ROI).
The test dataset (a $50\times50$ pixel scene rendered at quality
level~9 with antialiasing) was used as the input. Results show a
spill rate of \textbf{7.78\%} within the ROI --- the second highest
observed across all SPEC CPU2017 benchmarks evaluated in this study,
behind only \texttt{507.cactuBSSN\_r} (7.88\%) --- corresponding to
174 million detected spill events across 2.24 billion ROI instructions.
The ray-tracing pipeline simultaneously maintains ray state, surface
intersection records, shading context, and texture parameters, creating
extreme register pressure throughout the per-pixel rendering loop.
\end{abstract}

% =========================================================
\section{Benchmark Overview}

\texttt{511.povray\_r} is the SPEC CPU2017 rate variant of POV-Ray
3.6.1 (SPEC-modified), a ray-tracing renderer that produces
photorealistic images from scene descriptions. The benchmark exercises
recursive ray tracing with reflection, refraction, and shadow rays;
procedural texture evaluation; antialiasing with supersampling;
and bounding-box hierarchy traversal. Each rendered pixel requires
an unpredictable number of recursive ray casts, each evaluating
intersection tests against scene geometry, lighting models, and
multi-layered texture functions.

The RISC-V binary (\texttt{povray\_r.riscv}) was compiled with gem5
ROI markers inserted in \texttt{src/povray.cpp}, wrapping
\texttt{frontend.StartRender()} and the cooperate loop:

\begin{itemize}
    \item \texttt{m5\_work\_begin(0,0)} --- before \texttt{StartRender()}
    \item \texttt{m5\_work\_end(0,0)}   --- after the render completes
\end{itemize}

This ROI encompasses the complete per-pixel ray-tracing computation,
including all recursive ray spawning, shading, and texture evaluation.

\subsection{Input Datasets}

SPEC CPU2017 provides three dataset sizes for \texttt{povray\_r},
differing in image resolution, scene complexity, and rendering quality
(Table~\ref{tab:inputs}).

\begin{table}[H]
\centering
\caption{povray\_r Input Dataset Sizes}
\label{tab:inputs}
\begin{tabular}{lll}
\toprule
\textbf{Dataset} & \textbf{Resolution} & \textbf{Scene} \\
\midrule
test    & $50\times50$           & SPEC-benchmark-test.pov \\
train   & $\sim$$200\times150$   & SPEC-benchmark-train.pov \\
refrate & $1280\times768$        & SPEC-benchmark-ref.pov \\
\bottomrule
\end{tabular}
\end{table}

The \textbf{test} dataset was selected for this study. The input
\texttt{SPEC-benchmark-test.ini} specifies a $50\times50$ pixel
render at quality level~9 (maximum) with antialiasing
(\texttt{Antialias\_Depth=3}). Although the image is small, each
pixel requires the full recursive ray-tracing pipeline, including
bounding-box traversal, surface intersection, Phong shading, and
procedural texture evaluation using the POV-Ray standard include
libraries (\texttt{colors.inc}, \texttt{textures.inc},
\texttt{metals.inc}, \texttt{woods.inc}, \texttt{skies.inc}).

% =========================================================
\section{Simulation Setup}

\subsection{gem5 Configuration}

\begin{itemize}
    \item \textbf{CPU model:}   TimingSimpleCPU
    \item \textbf{ISA:}         RISC-V RV64GC
    \item \textbf{Caches:}      enabled (default L1 I/D + L2)
    \item \textbf{Memory:}      4\,GB simulated physical RAM
    \item \textbf{Spill mode:}  summary only (\texttt{spill\_verbose=False})
\end{itemize}

The full gem5 invocation was:

\begin{verbatim}
./build/RISCV/gem5.opt
  --outdir=benchmarks/povray/m5out_test
  configs/deprecated/example/se.py
  --cmd=benchmarks/povray/povray_r.riscv
  --options="benchmarks/povray/
             SPEC-benchmark-test.ini"
  --cpu-type=TimingSimpleCPU
  --caches
  --mem-size=4GB
\end{verbatim}

The ini file includes a \texttt{Library\_Path} directive pointing to
the benchmark directory so that POV-Ray's standard include files
(\texttt{*.inc}) are found at runtime.

% =========================================================
\section{Simulation Results}

The simulation completed successfully. POV-Ray rendered the
$50\times50$ pixel scene and wrote the output TGA image file.

\subsection{Pre-ROI Context}

Before the render ROI, the binary executed POV-Ray initialisation:
parsing the ini file, loading and compiling the scene description
(\texttt{.pov} file), parsing all \texttt{\#include} directives
(six standard include libraries), and building the bounding-box
hierarchy. Table~\ref{tab:preroi} shows the global counters at
ROI entry.

\begin{table}[H]
\centering
\caption{Global Counters at ROI Entry}
\label{tab:preroi}
\begin{tabular}{lr}
\toprule
\textbf{Metric} & \textbf{Value} \\
\midrule
Total instructions (pre-ROI) & 890,500 \\
Total spills detected (pre-ROI) & 48,951 \\
Pre-ROI spill rate & 5.50\% \\
\bottomrule
\end{tabular}
\end{table}

The pre-ROI spill rate of 5.50\% is the second highest pre-ROI rate
across all evaluated benchmarks, reflecting the complexity of
POV-Ray's scene compilation phase: parsing recursive \texttt{\#include}
files, constructing texture trees, and building the scene graph all
involve heavily recursive C++ call stacks.

\subsection{ROI Spill Statistics}

Table~\ref{tab:roi} presents the spill statistics collected within
the ROI.

\begin{table}[H]
\centering
\caption{ROI Spill Statistics --- 511.povray\_r (test, $50\times50$)}
\label{tab:roi}
\begin{tabular}{lr}
\toprule
\textbf{Metric} & \textbf{Value} \\
\midrule
ROI instructions         & 2,240,865,845 \\
ROI stores               &   287,983,968 \\
ROI loads                &   716,257,788 \\
\midrule
\textbf{ROI spills}      & \textbf{174,359,922} \\
\textbf{Spill rate}      & \textbf{7.7809\%} \\
\bottomrule
\end{tabular}
\end{table}

\subsection{Timing}

\begin{table}[H]
\centering
\caption{Simulation Timing}
\label{tab:timing}
\begin{tabular}{lr}
\toprule
\textbf{Metric} & \textbf{Value} \\
\midrule
Simulated time           & 4.97\,s \\
Host wall time           & 1{,}528.79\,s ($\approx$25.5\,min) \\
Total simulated instructions & 2,242,138,784 \\
\bottomrule
\end{tabular}
\end{table}

% =========================================================
\section{Analysis}

\subsection{Spill Rate Interpretation}

The measured spill rate of \textbf{7.78\%} is the second highest
across all SPEC CPU2017 benchmarks evaluated in this study
(Table~\ref{tab:compare}), behind only \texttt{507.cactuBSSN\_r}
(7.88\%) by a margin of just 0.10 percentage points. Approximately
1 in 13 instructions inside the ROI is a register spill.

\begin{table}[H]
\centering
\caption{Cross-Benchmark Spill Rate Comparison}
\label{tab:compare}
\begin{tabular}{lrr}
\toprule
\textbf{Benchmark} & \textbf{ROI Spills} & \textbf{Spill Rate} \\
\midrule
507.cactuBSSN\_r & 3,782,099,438 & 7.8756\% \\
511.povray\_r    &   174,359,922 & \textbf{7.7809\%} \\
500.perlbench\_r &     1,070,055 & 7.2838\% \\
502.gcc\_r       &     1,108,838 & 6.9185\% \\
526.blender\_r   &    61,845,869 & 6.2193\% \\
523.xalancbmk\_r &    18,381,250 & 4.6427\% \\
541.leela\_r     & 1,158,236,271 & 4.4438\% \\
519.lbm\_r       &   274,130,440 & 4.2130\% \\
531.deepsjeng\_r & 1,928,903,656 & 4.1749\% \\
544.nab\_r       &   151,310,616 & 2.2612\% \\
538.imagick\_r   &     2,424,899 & 2.0601\% \\
557.xz\_r        &    38,411,428 & 1.9069\% \\
505.mcf\_r       &   446,113,136 & 1.8263\% \\
508.namd\_r      &   327,718,955 & 1.0269\% \\
\bottomrule
\end{tabular}
\end{table}

\subsection{Store vs.\ Load Balance}

The load count (716.3M) is $2.49\times$ higher than the store count
(288.0M). This load-dominated ratio reflects the ray-tracing access
pattern: each intersection test reads the geometry of candidate
objects (bounding boxes, sphere parameters, triangle vertices), and
each shading step reads texture parameters and material properties,
while writes occur only when recording hit records and writing output
pixels. The many levels of recursive ray spawning amplify the read
dominance as each ray inherits the scene's read-only geometry data.

\subsection{Spill Fraction of Loads}

Within the ROI, the ratio of spills to total loads is:
\[
\frac{174{,}359{,}922}{716{,}257{,}788} \approx 24.3\%
\]
Approximately 1 in 4 load instructions is a spill reload. At each
level of the recursive ray-tracing call stack, POV-Ray must keep alive:
the ray origin and direction (6 floating-point values), the current
intersection record (hit distance, surface normal, UV coordinates),
the active lighting context, and the current texture evaluation state.
With RISC-V RV64GC providing 32 floating-point and 32 integer
registers, the compiler cannot keep all these values live across
nested function calls, resulting in pervasive spilling of both
integer pointers and floating-point ray parameters.

\subsection{Pre-ROI Spill Density}

The pre-ROI spill rate of 5.50\% ($48{,}951$ spills over $890{,}500$
instructions) is the second highest pre-ROI rate of all evaluated
benchmarks (behind only the 11.83\% of \texttt{cactuBSSN\_r}).
This high rate reflects POV-Ray's scene compilation phase, which
involves parsing deeply nested \texttt{\#include} files, recursively
evaluating texture macros, and constructing complex scene graph
objects — all highly recursive operations that produce severe
register pressure even before the first ray is cast.

\subsection{ROI Coverage}

The ROI accounts for $2{,}240{,}865{,}845$ out of $2{,}242{,}138{,}784$
total simulated instructions:
\[
\frac{2{,}240{,}865{,}845}{2{,}242{,}138{,}784} \approx 99.94\%
\]
This is among the highest ROI coverage fractions of all evaluated
benchmarks. The negligibly small pre-ROI phase (only 890,500
instructions) confirms that scene parsing is extremely fast relative
to the rendering computation, and the ROI captures nearly the
complete simulated workload.

\subsection{Ray-Tracing Register Pressure}

The 7.78\% spill rate is structurally driven by the recursive nature
of ray tracing. Each recursive call to cast a reflection, refraction,
or shadow ray requires saving the caller's full ray state to the stack.
At depth $d$, the call stack contains $d$ complete ray records
simultaneously. With POV-Ray's maximum recursion depth and RISC-V's
32 available registers (of which only $\sim$27 are usable for
general computation after accounting for the frame pointer, stack
pointer, and return address), even a modest recursion depth of 3--4
levels exceeds the register file capacity for the combined floating-point
and integer state. This forced spilling is compounded by the vectorised
texture evaluation functions, which also maintain large numbers of
simultaneously live floating-point intermediates. The result is a
spill pattern that closely resembles \texttt{cactuBSSN\_r}'s tensor
field operations: both workloads suffer from too many simultaneously
live floating-point values relative to the RISC-V register file size.

% =========================================================
\section{Conclusion}

The \texttt{511.povray\_r} benchmark exhibits a register spill rate
of 7.78\% within the rendering ROI on RISC-V RV64GC, the second
highest across all 14 evaluated SPEC CPU2017 benchmarks. With 174
million spill events across 2.24 billion ROI instructions and a
99.94\% ROI coverage, the benchmark is comprehensively characterised
by this run. The 2.49:1 load-to-store ratio reflects the read-heavy
geometry and texture access pattern of ray tracing, and the 24.3\%
spill-to-load fraction arises from the recursive ray state that must
be preserved across each spawned ray call. The near-identical spill
rate to \texttt{cactuBSSN\_r} (7.88\%) reveals that recursive
floating-point workloads and structured tensor-field computations
both approach the effective limit of the RISC-V RV64GC register
file for sustained floating-point computation.

\end{document}
